%% A realistic acmart submission, shaped like ACM's own samples/sigconf.tex.
%% Used to test venue-shift; not a real paper.
\documentclass[sigconf]{acmart}

\usepackage{booktabs}
\usepackage{hyperref}
\usepackage{pgfplots}

\setcopyright{acmlicensed}
\copyrightyear{2026}
\acmYear{2026}
\acmDOI{10.1145/3654321.3654322}
\acmConference[CHI '26]{CHI Conference on Human Factors in Computing
  Systems}{April 25--30, 2026}{Yokohama, Japan}
\acmISBN{978-1-4503-9999-9/26/04}

\begin{document}

\title{Trust Calibration in Highly Automated Driving}

\author{Mark Colley}
\authornote{Corresponding author.}
\orcid{0000-0001-5207-5029}
\email{mark.colley@uni-ulm.de}
\affiliation{%
  \institution{Institute of Media Informatics, Ulm University}
  \city{Ulm}
  \state{Baden-W\"urttemberg}
  \country{Germany}
}

\author{Jane Doe}
\email{jane.doe@example.edu}
\affiliation{%
  \institution{Example University}
  \city{Boston}
  \state{Massachusetts}
  \country{USA}
}

\begin{abstract}
  Automated vehicles must communicate their intent to build appropriate trust.
  We report a within-subjects study with 24 participants examining how takeover
  requests shape trust over repeated exposure.
\end{abstract}

\begin{CCSXML}
<ccs2012>
  <concept>
    <concept_id>10003120.10003121</concept_id>
    <concept_desc>Human-centered computing~Human computer interaction (HCI)</concept_desc>
    <concept_significance>500</concept_significance>
  </concept>
</ccs2012>
\end{CCSXML}
\ccsdesc[500]{Human-centered computing~Human computer interaction (HCI)}

\keywords{automated driving, trust, takeover request, user study}

\maketitle

\section{Introduction}
Trust in automated vehicles is neither uniformly good nor uniformly
bad~\cite{lee2004trust}. We examine how it develops across repeated takeovers.

\begin{figure*}[t]
  \centering
  \includegraphics[width=\textwidth]{figures/apparatus}
  \caption{The driving simulator used in the study.}
  \Description{A fixed-base driving simulator with three screens.}
  \label{fig:apparatus}
\end{figure*}

\section{Method}
We ran a within-subjects experiment with 24 participants.

\begin{table}[t]
  \caption{Participant demographics.}
  \label{tab:demographics}
  \begin{tabular}{lr}
    \toprule
    Measure & Value \\
    \midrule
    Participants & 24 \\
    Mean age & 27.4 \\
    \bottomrule
  \end{tabular}
\end{table}

\section{Results}
Trust increased over exposure, $F(2,46) = 8.1$, $p < .001$.

\section{Conclusion}
Repeated exposure calibrates trust.

\bibliographystyle{ACM-Reference-Format}
\bibliography{references}

\end{document}
